% Part:sets-functions-relations
% Chapter: size-of-sets
% Section: equinumerous-sets

\documentclass[../../../include/open-logic-section]{subfiles}

\begin{document}

\olfileid[pt]{sfr}{siz}{equ}

\olsection{Equinumerosidade}

Temos uma noção intuitiva do ``tamanho'' dos conjuntos, que funciona bem para
conjuntos finitos. E quanto aos infinitos? Se quisermos uma forma formal de
comparar os tamanhos de dois conjuntos de \emph{qualquer} cardinalidade, é boa
ideia começar por definir quando dois conjuntos têm o mesmo tamanho. Eis Frege:
\begin{quote}
  Se um garçom quer ter certeza de que pôs exatamente tantas facas
  quantos pratos sobre a mesa, não precisa contar nem uns nem outros, se simplesmente
  puser uma faca à direita de cada prato, de modo que cada faca na mesa fique à
  direita de algum prato. Os pratos e as facas ficam assim correlacionados de
  modo único, e de fato por essa mesma relação espacial. \citep[§~70]{Frege1884}
\end{quote}
A intuição central desta passagem traduz-se numa definição formal:

\begin{defn}\ollabel{comparisondef}
  $A$ é \emph{equinumeroso} a $B$, escrevendo-se $\cardeq{A}{B}$, sse existir
  !!a{bijection} $f \colon A \to B$.
\end{defn}

\begin{prop}\ollabel{equinumerosityisequi}
A relação de equinumerosidade é uma relação de equivalência.
\end{prop}

\begin{proof} 
É preciso mostrar que a equinumerosidade é reflexiva, simétrica e transitiva. Sejam
$A, B$ e $C$ conjuntos.

\emph{Reflexividade.} A aplicação identidade $\Id{A} \colon A \to A$, com
$\Id{A} (x) = x$ para todo $x \in A$, é !!a{bijection}. Logo $\cardeq{A}{A}$.

\emph{Simetria.} Suponhamos $\cardeq{A}{B}$, isto é, que existe !!a{bijection}
$f\colon A \to B$. Como $f$ é !!{bijective}, a inversa $f^{-1}$ existe e também é
!!{bijective}. Logo $f^{-1}\colon B \to A$ é !!a{bijection}, pelo que
$\cardeq{B}{A}$.

\emph{Transitividade.} Suponhamos $\cardeq{A}{B}$ e $\cardeq{B}{C}$, isto é, que
existem !!{bijection}s $f\colon A \to B$ e $g\colon B \to C$. Então a composição
$\comp{f}{g}\colon A \to C$ é !!{bijective}, logo $\cardeq{A}{C}$.
\end{proof}

\begin{prop}
Se $\cardeq{A}{B}$, então $A$ é !!{enumerable} se e só se $B$ o é.
\end{prop}

\begin{editorial}
A prova seguinte usa \olref[enm]{defn:enumerable} se \olref[enm]{sec} estiver
incluída e \olref[enm-alt]{defn:enumerable} caso contrário.
\end{editorial}

\begin{proof}
Suponhamos $\cardeq{A}{B}$, logo existe alguma !!{bijection} $f \colon A \to B$, e
suponhamos que $A$ é !!{enumerable}.
\oliflabeldef{sfr:siz:enm:defn:enumerable}{
  Então ou $A = \emptyset$ ou existe uma função !!{surjective} $g\colon \PosInt \to A$.
  Se $A = \emptyset$, então $B = \emptyset$ também (caso contrário existiria
  !!a{element}~$y \in B$ sem $x \in A$ com $g(x) = y$). Se, por outro lado,
  $g\colon \PosInt \to A$ for !!{surjective}, então $\comp{g}{f} \colon \PosInt \to
  B$ é !!{surjective}. De fato, seja $y \in B$. Como $f$ é !!{surjective}, existe
  $x \in A$ tal que $f(x) = y$. Como $g$ é !!{surjective}, existe $n \in \PosInt$
  tal que $g(n) = x$. Logo,
\[
(\comp{g}{f})(n) = f(g(n)) = f(x) = y
\]
e portanto $\comp{g}{f}$ é !!{surjective}. Temos que $\comp{g}{f}$ é uma enumeração
de~$B$, logo $B$~é !!{enumerable}.}
{
  Então ou $A  = \emptyset$ ou existe !!a{bijection}~$g$ cujo contradomínio é $A$ e
  cujo domínio é $\Nat$ ou um segmento inicial de números naturais. Se $A =
  \emptyset$, então $B = \emptyset$ também (caso contrário existiria algum~$y \in B$
  sem $x \in A$ tal que $g(x) = y$). Suponhamos então que temos a
  !!{bijection}~$g$. Então $\comp{g}{f}$ é !!a{bijection} com contradomínio~$B$ e
  domínio o mesmo que o de~$g$ (isto é, $\Nat$ ou um segmento inicial), logo $B$ é
  !!{enumerable}.}

Se $B$ é !!{enumerable}, obtemos que $A$ é !!{enumerable} repetindo o argumento com
a !!{bijection} $f^{-1}\colon B \to A$ em lugar de~$f$.
\end{proof}

\begin{prob}
Mostre que, se $\cardeq{A}{C}$ e $\cardeq{B}{D}$, e $A \cap B = C \cap D =
\emptyset$, então $\cardeq{A \cup B}{C \cup D}$.
\end{prob}

\begin{prob}
Mostre que, se $A$ é infinito e !!{enumerable}, então $\cardeq{A}{\Nat}$.
\end{prob}

\end{document}
